\documentclass[../../../include/open-logic-section]{subfiles}

\begin{document}

\olfileid{sth}{cardinals}{cp}
\olsection{Cantor原则}

让我们回想一下 \olref[ordinals][vn]{sec}。当时我们在讨论良序集，并提出希望能有某种对象来充当良序的\emph{代表}。基于这一想法，我们引入了序数，并在
\olref[ordinals][ordtype]{ordtypesworklikeyouwant} 中证明了它们的表现确实符合我们的预期，即：
\[
	\ordtype{A, <} = \ordtype{B, \lessdot} 
	\text{ 当且仅当 } \tuple{A, <} \isomorphic \tuple{B, \lessdot}.
\]
再往前回想 \olref[sfr][siz][equ]{sec}。在那里，我们朴素地引入了集合“大小”的概念。具体而言，我们说两个集合是等势的，记作 $\cardeq{A}{B}$，当且仅当存在一个双射 $f \colon A \to
B$。这个概念本质上比良序的概念\emph{更简单}：我们只关心是否存在双射，而不关心双射是否“保持了某种结构”（序同构则要求这一点）。

由此自然产生了一个显而易见的想法。正如我们引入某些对象——即\emph{序数}——来标度良序一样，我们也可以引入某些对象——即\emph{基数}——来标度大小。这正是本章的目标。

在说明基数究竟是什么之前，我们应当先确立它们理应满足的一条原则。用 $\card{X}$ 表示集合 $X$ 的基数，我们希望基数满足：
\[
	\card{A} = \card{B} \text{ 当且仅当 } \cardeq{A}{B}.
\]
我们将此称为\emph{康托尔}原则，因为康托尔很可能是第一个对此有清晰认识的人。（关于它与\emph{休谟}原则的关系，我们将在 \olref[hp]{sec} 中作进一步讨论。）因此，我们的目标就是为每个 $X$ 定义 $\card{X}$，使之满足康托尔原则。

\end{document}